\bookchapter{Kapitel 35: Kinder und Smartphones – Erziehungsleitfaden fuer das digitale Zeitalter}

\chapterintro{Dieses Kapitel behauptet nicht, das Handy sei »ein Ungeheuer«, noch dass »je frueher desto klueger«. Wir gehen nach Altersstufen vor: Geraete bereitstellen, Regeln setzen, Medienkompetenz vermitteln, Gesundheit schuetzen – Schritt fuer Schritt.}

\booksection{Wann gibt man dem Kind ein Handy?}

Es gibt keine Standardantwort, aber altersbezogene Richtwerte. Zu frueh: Der praefrontale Kortex ist noch nicht ausgereift, die Selbstkontrolle kann mit dem »Verlockungs-Design« der Algorithmen nicht mithalten. Zu spaet: Moegliche soziale Isolation unter Gleichaltrigen. Hier die Empfehlungen nach Alter:

\begin{itemize}
  \item \textbf{3–6 Jahre}: Möglichst nicht ans Handy lassen. Bei Nutzung: höchstens 15 Minuten pro Sitzung, Inhalte sollten von Eltern ausgewählte pädagogische Animationen sein, und es muss\textbf{gemeinsam ansehen}. In dieser Phase hat der Bildschirm den groessten Einfluss auf Sehkraft und Aufmerksamkeitsentwicklung.
  \item \textbf{7–9 Jahre}: Kann auf dem Gerät der Eltern "geliehen" werden – für Online-Unterricht, Geschichten und Nachschlagen. Die Regel "vorher fragen, danach zurückgeben" aufbauen. Kein eigenes Handy.
  \item \textbf{10–12 Jahre}: Ein "Feature Phone" oder eine Kinderuhr kann den Schulweg-Kontakt lösen. Wenn ein Smartphone nötig ist, strenge Elternkontrolle aktivieren.
  \item \textbf{Ab 13 Jahren}: In der Mittelstufe wächst das soziale Bedürfnis – ein Smartphone ist möglich, aber mit einem "Nutzungsvertrag" und regelmäßiger Reflexion.
\end{itemize}

\begin{tcbinfo}{Das Prinzip »So viel wie noetig«}
Das erste Handy des Kindes muss kein Flaggschiff sein. Ueberschuessige Prozessorleistung wird oft fuer Spiele genutzt. Ein Mittelklasse-Geraet mit grossem Akku und augenfreundlichem Display ist leichter zu kontrollieren. Budget-Richtwert: 1500–2500 Yuan genuegen fuer die meisten Schueler.

\end{tcbinfo}

\booksubsection{Smartwatch als uebergangsloesung}

Wenn es nur darum geht, nach der Schule erreichbar zu sein, ist eine Kinderuhr besser als ein Handy. Sie kann nur telefonieren, Sprachnachrichten senden und Standort melden – kein Browser, keine Kurzvideos, von Natur aus eingeschraenkt:

\begin{itemize}
  \item \textbf{Xiao Tian Cai Telefonuhr}: Vollständigstes Ökosystem. "Antippen zum Hinzufügen" ist der Kernverkaufspunkt unter Gleichaltrigen, aber die App-Ökosphäre ist relativ geschlossen.
  \item \textbf{Huawei Kinderuhr}: Gute Integration mit Huawei-Handys, präzise Ortung, Videotelefonie, gutes Preis-Leistungs-Verhältnis.
  \item \textbf{Xiaomi Kinderuhr}: Günstig, Grundfunktionen vollständig, geeignet für Familien mit begrenztem Budget.
\end{itemize}

Beim Uhrenkauf darauf achten: Ein Modell mit\textbf{4G All-Net}waehlen (2G/3G werden schrittweise abgeschaltet) und die Tarife fuer Kinder-Pakete beim Anbieter pruefen.

\booksection{Elternkontrollen – der komplette Ueberblick}

Bevor das Geraet dem Kind uebergeben wird, ist die erste Massnahme nicht das Einschalten, sondern das Einrichten der Elternkontrolle. Beide grossen Plattformen bieten native Loesungen:

\booksubsection{iOS: Bildschirmzeit}

\begin{itemize}
  \item \textbf{App-Limits}: Einstellungen > Bildschirmzeit > App-Limits. Für jede App-Kategorie kann ein tägliches Limit gesetzt werden (z. B. Spiele 30 Minuten).
  \item \textbf{Sperrzeit}: Tägliche "Licht-aus-Zeit" festlegen (z. B. 22:00–7:00). In dieser Zeit werden alle Apps grau und deaktiviert, nur Telefon ist verfügbar.
  \item \textbf{Inhaltsbeschraenkungen}: Filterung von Filmen, Webseiten und App-Installationen nach Altersgruppe sowie Sperrung von In-App-Käufen.
  \item \textbf{Familienfreigabe}: Verwaltung des Kindergeräts per elterlichem iPhone aus der Ferne – Download-Anforderungen genehmigen, Nutzungsberichte ansehen.
\end{itemize}

\booksubsection{Android: Google Family Link / Marken-Elternkontrolle}

\begin{itemize}
  \item \textbf{Google Family Link}: Googles offizielle Lösung – Geräte-Limits aus der Ferne, Gerät sperren, Aktivitätsberichte ansehen. In China je nach Modell verfügbar.
  \item \textbf{Xiaomi Kinder-Bereich}: Einstellungen > Kinder-Bereich. Nutzungsdauer und Whitelist-Apps festlegen; zum Verlassen Eltern-Passwort nötig.
  \item \textbf{Huawei Schüler-Modus}: Einstellungen > Gesunde Handy-Nutzung > Schüler-Modus. Nutzungszeiten und Apps einschränken, Fernverwaltung unterstützt.
  \item \textbf{OPPO/vivo Kinder-Modus}: In "Einstellungen" oder der dedizierten "Kinder-Bereich"-App aktivieren. Ähnliches Prinzip.
\end{itemize}

\begin{center}
\begin{tabularx}{\linewidth}{|X|X|X|}
\hline
\textbf{Funktion} & \textbf{iOS Bildschirmzeit} & \textbf{Android Family Link / Marken-Lösung} \tabularnewline
\hline
App-Limits & Unterstützt, nach Kategorie & Unterstützt, nach App \tabularnewline
Sperrzeit & Unterstützt & Unterstützt (Marken-Lösungen variieren) \tabularnewline
Inhaltsfilter & Auf Systemebene, gründlicher & Verändert je nach Marke, uneinheitlich \tabularnewline
Fernverwaltung & Familienfreigabe & Family Link / Marken-Cloud \tabularnewline
App-Installationsfreigabe & Elterliche Freigabe erforderlich & Family Link unterstützt \tabularnewline
\hline
\end{tabularx}
\end{center}

\begin{tcbwarn}{Elternkontrolle ist keine "einmalige Lösung"}
Kinder finden immer Wege, sie zu umgehen (im Internet gibt es massenhaft Anleitungen zum Knacken der Bildschirmzeit). Technik ist nur die Grundlinie; die echte Verteidigungslinie sind Kommunikation und Vertrauen zwischen Eltern und Kind. Empfehlung: Den Nutzungsbericht jeden Monat zusammen mit dem Kind besprechen, nicht nur einseitig überwachen.

\end{tcbwarn}

\booksection{III. Jugendmodus}

Alle großen Content-Plattformen haben einen integrierten “Jugendmodus”, der nach Aktivierung Nutzungsdauer begrenzt, schädliche Inhalte filtert und Aufladungen/Trinkgelder sperrt. Dies ist gesetzlich vorgeschrieben, aber\textbf{muss von den Eltern aktiv aktiviert und mit Passwort geschützt werden}:

\begin{itemize}
  \item \textbf{Douyin/Kuaishou Jugendmodus}: Täglich 40 Minuten (22:00–6:00 gesperrt), nur ausgewählte populärwissenschaftliche und pädagogische Videos, kein Live-Streaming, kein Trinkgeld, keine Aufladungen.
  \item \textbf{WeChat Jugendmodus}: WeChat Pay einschränken (Eltern müssen Limit freigeben), öffentliche Accounts und Video-Kanal einschränken, "Schütteln" und "In der Nähe" standardmäßig deaktiviert.
  \item \textbf{Anti-Sucht-System für Spiele}: Nationale Echt Namens-Verifizierung + Zeitbeschränkung. Minderjährige können nur freitags, samstags, sonntags und an Feiertagen von 20:00–21:00 Uhr spielen, täglich 1 Stunde.
\end{itemize}

\textbf{Bei der Echt Namens-Verifizierung nicht den Ausweis der Eltern “leihen”}\par

Viele Kinder nutzen den Ausweis der Eltern, um Spielekonten zu registrieren und die Anti-Sucht-Sperre zu umgehen. Das verstößt gegen die Nutzungsbedingungen und nimmt dem Kind den Schutz. Eltern sollten ihre Zahlungs-PINs und Ausweisdaten sicher aufbewahren und nicht aus Bequemlichkeit “leihen”.

\booksection{IV. Medienkompetenz-Erziehung}

Technik blockt 90 \% der Risiken, die restlichen 10 \% hängen vom Urteilsvermögen des Kindes ab. Statt “Internetverbot” ist es wichtiger, ihnen beizubringen, “wie man das Internet nutzt”:

\begin{itemize}
  \item \textbf{: Die erweiterte Garderobe, etwas langsamer, geeignet für Filme, Fotos und andere große Dateien. Achtung: Viele Flagship-Smartphones unterstützen keine SD-Karten mehr.}: Fremden nicht den echten Namen, die Schule oder die Adresse verraten; keine persönlichen Fotos, Ausweise oder Familienbilder senden; Standortfunktion standardmäßig aus, nur bei Bedarf an.
  \item \textbf{Cybermobbing erkennen und reagieren}: Beschimpfungen, Ausgrenzung und bösartige Fotomanipulation sind Mobbing. Sagen Sie dem Kind: Nicht zurückschlagen, Beweise nicht löschen, sofort den Eltern sagen. In WeChat kann man Nachrichten lange drücken, um sie zu melden – die Plattform geht gegen den Täter vor.
  \item \textbf{Kritischer Umgang mit Inhalten}: Bringen Sie dem Kind bei, "Fakten" von "Meinungen" zu unterscheiden, Clickbait und Sensationsmeldungen zu erkennen; zu verstehen, dass der Algorithmus nur empfiehlt, was man sehen will – nicht die ganze Welt; bei "Geschenke vom Himmel"-Werbung immer misstrauisch zu sein.
  \item \textbf{Netzwerksicherheit}: Keine fremden Links klicken, keine unbekannten QR-Codes scannen, keine Apps außerhalb offizieller Stores laden; Passwörter niemandem geben (auch nicht "besten Freunden"); wer nach Codes fragt, ist immer ein Betrüger.
\end{itemize}

\begin{tcbpractice}{Praxis: Zusammen mit dem Kind einen "Handy-Nutzungsvertrag" unterzeichnen}
Regeln aufschreiben, beide unterschreiben und gut sichtbar aufhängen – das wirkt besser als mündliche Absprachen. Empfohlene Inhalte:

\begin{itemize}
  \item Nutzungszeit: z. B. "Nach den Hausarbeiten 19:30–20:30".
  \item Sperrzonen: Am Esstisch, im Schlafzimmer vor dem Einschlafen, im Unterricht – Handy weg.
  \item Inhaltsgrenzen: Keine Apps installieren, ohne die Eltern zu informieren; keine Privatchats mit Fremden.
  \item Folgen bei Verstößen: z. B. "Bei einmaligem Überziehen wird die Nutzungsdauer am nächsten Tag halbiert".
  \item Elternversprechen: Nicht länger als vereinbart vor dem Kind "scrollen" – Vorbild sein.
\end{itemize}

\end{tcbpractice}

\booksection{V. Gesunde Nutzungsgewohnheiten}

Augen, Nacken und Schlaf von Kindern sind in der Entwicklung – der Handy-Schaden ist größer als bei Erwachsenen:

\begin{itemize}
  \item \textbf{Bildschirmzeit-Empfehlungen}: 3–6 Jahre: täglich max. 1 Stunde; 6–12 Jahre: Lernen unbegrenzt, Unterhaltung max. 1 Stunde; über 12 Jahre: mit dem Kind vereinbaren, Unterhaltung aber max. 2 Stunden.
  \item \textbf{Augenschonende Gewohnheiten}: Die "20-20-20-Regel" einhalten – alle 20 Minuten Bildschirm 20 Sekunden auf 20 Fuß (ca. 6 m) in die Ferne blicken; 1 Stunde vor dem Schlafengehen nicht aufs Handy, da Blaues Licht Melatonin hemmt und den Schlaf stört.
  \item \textbf{Haltungsmanagement}: Handy-Bildschirm auf Augenhöhe, nicht mit gesenktem Kopf und krummem Rücken; nach 30 Minuten aufstehen und bewegen.
  \item \textbf{Gemeinsame Familienregeln}: Esstisch, Schlafzimmer, Arbeitszimmer als "handyfreie Zonen"; beim gemeinsamen Filmschauen oder Outdoor-Aktivitäten alle Handys einsammeln – "Handy weg" als Familienkultur, nicht als Strafe für das Kind.
\end{itemize}

\begin{tcbinfo}{Eltern sind die beste "Firewall"}
Studien zeigen: Die Handy-Nutzungsdauer der Eltern korreliert stark mit der der Kinder. Wenn man vom Kind verlangt, nicht aufs Handy zu schauen, selbst aber bis spät in die Nacht Kurzvideos scrollt, werden die Regeln wirkungslos. Um das Kind zu ändern, zuerst das eigene Handy weglegen.

\end{tcbinfo}

\booksection{VI. Digitales Erbe und Kinder}

Dieser Abschnitt wird oft übersehen, ist aber wichtig. Die digitalen Vermögenswerte von Kindern (Accounts, Fotos, Guthaben, Spielgeräte) haben ebenfalls Erb- und Privatsphäre-Fragen:

\begin{itemize}
  \item \textbf{Vererbung digitaler Kinder-Vermögenswerte}: Wenn ein Kind verstirbt, ist der Ablauf für die Übertragung von Social-Media-Accounts, Spielkonten und Cloud-Fotos derselbe wie bei Erwachsenen – Sterbeurkunde und Verwandtschaftsnachweis erforderlich. Eltern sollten für wichtige Konten des Kindes von Anfang an einen "Legacy Contact" oder "Notfallkontakt" einrichten (z. B. Apple Legacy Contact, WeChat-Account-Vererbung).
  \item \textbf{Digitale Privatsphäre Minderjähriger}: Nach dem Gesetz zum Schutz persönlicher Daten ist für die Verarbeitung von Daten Minderjähriger unter 14 Jahren die Zustimmung des Vormunds erforderlich, und es müssen spezielle Datenschutzregeln gelten. Eltern sollten regelmäßig die Datenschutzrichtlinien der vom Kind genutzten Apps prüfen und "personalisierte Empfehlungen" und "Datenfreigabe" deaktivieren.
  \item \textbf{Empfehlungen zur Kontoregistrierung}: Für das Kind registrierte Accounts sollten eine von den Eltern kontrollierte E-Mail/Handynummer verwenden, um die Verwaltung in der Minderjährigkeit zu erleichtern; nach Volljährigkeit kann das Konto übergeben werden.
\end{itemize}

\begin{tcbwarn}{Die andere Seite des "digitalen Kinder-Erbes"}
Kinder werden von Geburt an “aufgezeichnet” – die von Eltern geposteten Wachstumsfotos, Videos der Kita, Gesundheitsdaten der Smartwatch. Diese Daten lassen sich nach dem Upload kaum vollständig löschen. Vor dem Posten überlegen: Wird das Kind in zehn Jahren dieses Foto peinlich finden? Die digitale Persönlichkeit des Kindes respektieren – beginnen Sie damit, das “Kinder-Posting” zurückzuschrauben.

\end{tcbwarn}

\begin{tcbsummary}{Kapitelzusammenfassung}
\begin{itemize}
  \item Es gibt keine Standardantwort für das Alter der Handyvergabe, aber 3–6 Jahre möglichst gar nicht; vor 13 Jahren mit Kinderuhr oder Feature Phone überbrücken
  \item "Reicht aus" ist gut: Das erste Handy sollte ein Einsteigermodell sein, kein Flaggschiff – zu viel Leistung bedeutet zu viel Versuchung
  \item iOS-Bildschirmzeit und Android Family Link / Marken-Elternkontrolle sind zwei technische Verteidigungslinien – in Kombination sicherer
  \item Jugendmodus und Anti-Sucht-System sind gesetzlich vorgeschrieben, müssen aber von Eltern aktiv aktiviert und mit Passwort geschützt werden
  \item Medienkompetenz ist wichtiger als Technik: Kinder lehren, Mobbing zu erkennen, Privatsphäre zu schützen und Inhalte kritisch zu hinterfragen
  \item Gesunde Gewohnheiten beginnen bei den Eltern – erst wenn Sie das Handy weglegen, legt das Kind seines weg
  \item Digitales Kinder-Erbe und Privatsphäre sind neue Themen – "Legacy Contact" einrichten und Kinder-Postings zurückschrauben
\end{itemize}

\end{tcbsummary}
